\documentclass[11pt]{article}
\usepackage{amsmath,amssymb,amsthm,hyperref}
\newtheorem{theorem}{Theorem}
\newtheorem{lemma}[theorem]{Lemma}
\newtheorem{corollary}[theorem]{Corollary}
\newcommand{\Q}{\mathbb Q}
\newcommand{\CH}{\operatorname{CH}}
\newcommand{\cl}{\operatorname{cl}^{\rm can}}
\title{Closing the curvature realization test: projective retractions and boundary sources}
\author{Hand mathematics before Lean; author self-review}
\date{3 November 2026}
\begin{document}\maketitle

\section{Exact input and the distinction being tested}
Start at commit \texttt{6c78f4d57fcad0e2309cc0bba3417952a58f4d5b}.
The input is every connected smooth projective complex $n$-fold $X$,
$1\le p\le n$, and a prescribed rational Hodge class
$\alpha\in H^{2p}(X,\Q)\cap H^{p,p}(X)$.
The required final output remains a cycle $Z_\alpha\in\CH^p(X)_\Q$
with $\cl_X Z_\alpha=\alpha$. Connected components are treated separately.

The central-torus curvature/opposedness theorem and finite homogeneous
curvature normal form are established in their stated categories in the
earlier hand derivations. Their Lean counterparts are uncompiled drafts.
The torsion principal-torus projective construction is also established.
The question here is the remaining proposed passage from the higher
prescribed-class nullhomotopy to actual projective geometry. It is tested
without assuming a cycle for $\alpha$.

Three routes are executed: an arbitrary smooth projective realization of
the target-dependent homotopy fiber, a realization after finitely many
projective folds/alterations, and an open realization with a high-codimension
algebraic boundary. The first two have a uniform obstruction. The third
has an explicit source formula and an exact boundary threshold, but no
universal existence theorem for the required boundary is asserted.

\section{Construct the return for every smooth projective surjective fold}
\begin{theorem}[Actual projective return, without a torus hypothesis]\label{thm:return}
Let $f:Y\to X$ be any surjective morphism of connected smooth projective
complex varieties, $\dim Y=n+r$. Choose a very ample line $H$ on $Y$ and
put $h=c_1^{\CH}(H)$. There is a positive integer $d$ such that
\[
 f_*(h^r)=d[X]\quad\hbox{in }\CH^0(X).
\]
The actual rational Chow operation and its canonical Betti comparison are
\[
 R_{\CH}(B)=d^{-1}f_*(h^r B),\qquad
 R_B(b)=d^{-1}f_*(\cl(h)^r b).
\]
Both preserve codimension (respectively total cohomological degree), and
\[
 R_{\CH}f^*=\mathrm{id},\qquad R_B f^*=\mathrm{id},\qquad
 \cl_X R_{\CH}=R_B\cl_Y.                         \tag{1}
\]
No smoothness, flatness, connectedness or constant dimension of the fibers
is required. In particular $f^*$ is injective in every rational Betti degree.
\end{theorem}
\begin{proof}
The generic fiber has dimension $r$. Its intersection with $r$ general
hyperplanes from $H$ has positive finite degree $d$. Equivalently, intersect
$Y$ with those hyperplanes; its dominant components give an actual effective
multisection cycle of degree $d$, while components with smaller-dimensional
image push to zero in dimension $n$. Thus $f_*(h^r)=d[X]$.
For $r=0$ this is the positive generic degree of $f$.
The general pullback on Chow rings is available because both varieties
are smooth; it need not be a flat pullback. Its projection formula gives
$f_*(h^r f^*B)=f_*(h^r)B=dB$. This proves the first identity.
The canonical cycle map respects line Chern classes, products, proper
pushforwards and this pullback. The topological projection formula and
the comparison for $h^r$ give the Betti identity for every degree, not just
algebraic classes. Dividing by the positive $d$ is legitimate over $\Q$.
\end{proof}
For the exact Chow input used above see the Stacks Project,
Lemma 43.27.1, tag \href{https://stacks.math.columbia.edu/tag/0B2X}{0B2X},
checked 3 November 2026: its hypotheses are nonsingular projective varieties
and an arbitrary morphism, and its part (3) is precisely the required
projection formula. The choice of $H$ and the return need not be canonical;
the cycle-class map in (1) is canonical.

\begin{corollary}[No hidden disappearance under projective operations]\label{cor:defect}
For every such $f$ and prescribed $\alpha$,
\[
 \alpha\in\operatorname{im}\cl_X
 \quad\Longleftrightarrow\quad
 f^*\alpha\in\operatorname{im}\cl_Y.                \tag{2}
\]
Moreover the rational Hodge/cycle-image quotient for $X$ is a split
summand of that for $Y$. The same holds under any finite composable
sequence of these projective maps, using the composed returns.
\end{corollary}
\begin{proof}
Pull back an actual cycle for the forward implication. For the reverse,
apply $R_{\CH}$ to an actual cycle for $f^*\alpha$ and use (1).
The maps in (1) preserve Hodge type after the displayed product and
degree shift and preserve canonical cycle images. Hence they descend
to the quotients and still compose to identity. Composing split injections
and their reverse-order retractions proves the last assertion.
\end{proof}
This does not say a projective source on $Y$ cannot help construct a source
on $X$. It says that restoring projectivity or preserving Hodge type alone
has not supplied that source: its image must still be computed.

\section{Execute the prescribed-alpha higher-curvature construction}
Represent $\alpha$ by a classifying map $a_\alpha:X\to K(\Q,2p)$ and form
its actual homotopy fiber $t:T_\alpha\to X$. The chosen nullhomotopy gives
\[
 t^*\alpha=0.                                      \tag{3}
\]
This construction depends on the prescribed $\alpha$, not on a fixed
hyperplane class. Its rational cochain description adjoins a generator
$u$ of degree $2p-1$ with $du=\eta_\alpha$. One may also adjoin the
flat companion used by the paired GST model; (3) is unchanged.

\begin{theorem}[The direct projective realization attempt is impossible]\label{thm:noproper}
If $\alpha\ne0$, there is no surjective smooth-projective algebraic model
$f:Y\to X$ of $t:T_\alpha\to X$, even at the level of rational cohomology
compatible with the marked pullback from $X$. This remains true if finitely
many projective surjections, blowups between smooth varieties, or smooth
projective generically finite alterations over $X$ are inserted.
\end{theorem}
\begin{proof}
Compatibility would give $f^*\alpha=0$ by (3). But (1) would then give
$\alpha=R_B f^*\alpha=0$. The same contradiction applies to a composite.
At the cochain level it is the already established GST marked-base
boundary invariant: a retraction onto the closed base cannot kill a
nonzero base class which is a boundary. The new point is that actual
projective geometry constructs that retraction in all degrees, through
the explicit multisection above.
\end{proof}
In particular the theorem applies even when $\alpha$ already is algebraic.
It is an obstruction to this \emph{nullifying projective model}, not an
obstruction to algebraicity or a Hodge-conjecture counterexample. For
$p=1$ the punctured line bundle does kill $c_1(L)$, but it is open, not
projective. The compact elliptic torsor with nonzero rational Chern class
is likewise outside the smooth-projective category. These facts agree
with the proved curvature theorem and do not contradict it.

\section{Execute the boundary alternative and construct its actual cycle}
\begin{lemma}[Exact top support threshold]\label{lem:support}
Let $Y$ be smooth projective of complex dimension $m$ and $D\subset Y$
a closed algebraic subset of complex codimension at least $p\ge1$.
Write $D_1,\ldots,D_s$ for its irreducible codimension-$p$ components,
and $j:U=Y\setminus D\hookrightarrow Y$. Then
\[
 \ker\bigl(j^*:H^{2p}(Y,\Q)\to H^{2p}(U,\Q)\bigr)
 =\operatorname{span}_\Q\{\cl_Y[D_i]:1\le i\le s\}.       \tag{4}
\]
In particular the kernel is zero if $\operatorname{codim}_Y D>p$.
\end{lemma}
\begin{proof}
The cohomology-with-supports exact sequence identifies the kernel with
the image of $H_D^{2p}(Y,\Q)$. Remove the singular loci and pairwise
intersections of the $D_i$, and all higher-codimension components. The
removed algebraic set has complex codimension at least $p+1$ in $Y$.
Support cohomology below degree $2p+2$ on that set vanishes. To justify
the vanishing without a smoothness assumption on it, take a finite smooth
algebraic stratification, use the tubular-neighborhood Thom isomorphism
for each stratum (real codimension at least $2p+2$), and the exact
support sequences inductively. Consequently that removal leaves
$H_D^{2p}$ unchanged. The remaining $D_i$ are disjoint smooth connected
dense opens. Their complex normal orientations give
$H_D^{2p}\cong\bigoplus_i H^0(D_i^{\rm good},\Q)(-p)
\cong\Q(-p)^s$. The maps to $Y$ send these Thom generators to exactly
the canonical fundamental classes $\cl_Y[D_i]$.
There may be relations in their images; no independence of those images
or uniqueness of coefficients is asserted.
\end{proof}

\begin{theorem}[Actual boundary-to-original-X source formula]\label{thm:boundary}
Let $f:Y\to X$ and $H,r,d$ be as in Theorem~\ref{thm:return}.
Suppose an actual closed algebraic $D\subset Y$ of codimension at least
$p$ has been constructed and $j^*f^*\alpha=0$. Then rational coefficients
$b_i$ exist with $f^*\alpha=\sum_i b_i\cl_Y[D_i]$, and the actual cycle
\[
 \boxed{Z_\alpha={1\over d}\sum_i b_i\,
                f_*(h^r\cdot[D_i])\in\CH^p(X)_\Q}       \tag{5}
\]
has $\cl_X Z_\alpha=\alpha$. The coefficients are obtainable from a
finite rational linear system in the actual Betti classes of the boundary
components. Their nonuniqueness does not change the returned Betti class;
it is not asserted to leave the Chow class unchanged.
If $\alpha\ne0$, such a $D$ must have a codimension-$p$ component.
\end{theorem}
\begin{proof}
Use (4) on $f^*\alpha$. Then apply $R_B$ from (1) and canonical
comparison to the displayed finite rational cycle combination. Multiplying
by $h^r$ raises codimension to $p+r$; proper pushforward lowers it by $r$.
If no $D_i$ exists, (4) and the injectivity in (1) force $\alpha=0$.
\end{proof}

This is a completed source extraction from \emph{actual} geometric boundary
data, including singular or reducible $D$. It does not claim to have produced
that data for every input. In fact its existence condition is exactly
equivalent to $\alpha$ being algebraic: (5) proves one direction; for
the other, take $Y=X$, $f=\mathrm{id}$ and $D$ the union of the supports
of a rational cycle for $\alpha$. Thus the condition cannot be inserted
as an independent GST source-birth law. This is consistent with the
branch's exact Gysin-support/Chow-image theorem.

The threshold is sharp. For $X=\mathbf P^n$, $\alpha=h^p$ and a linear
subspace $D=\mathbf P^{n-p}$, projection from $D$ makes
$X\setminus D$ an affine-space bundle over $\mathbf P^{p-1}$.
Consequently $h^p$ restricts to zero and the codimension-$p$ boundary
returns its actual class. A boundary of codimension $p+1$ could not do
so by (4). This is a normalization check of (4), not a replacement of
the arbitrary-$X$ research target by a special-family program.

An ordinary smooth compactification of an open variety has a boundary
divisor. For $p>1$ this has codimension $1<p$, so (4) cannot be applied
by discarding its lower-codimension pieces. Their support classes involve
higher cohomology of the boundary; those are not automatically fundamental
cycle classes. This identifies the exact unsolved geometric implication
in the open continuation rather than declaring all nullhomotopies algebraic.

\section{Audit and checkpoint}
The curvature theorem survives the audit at its stated family scope.
The general projective retraction, the projective nullifier obstruction,
and the boundary extraction formula are proved here at hand-proof level,
by author self-review. New Lean drafts encode their algebraic consequences,
not projective varieties, topology, purity or a canonical cycle map instance.

The prescribed-input attempt genuinely executed was $T_\alpha$: it has
the computed marked class (3), and its direct projective-realization route
is ruled out by (1), including all finite projective compositions. The
open route has the explicit positive source formula (5). No universal
algebraic boundary realizing (3) is constructed. Its lower-codimension
boundary continuation remains open and requires new geometry; the support
equivalence is not counted as source-existence progress. The original
all-$X$ canonical Chow theorem therefore remains unproved.
\end{document}
